% ===================================================================
%  TABULKA č. 12 — Nádraží. — Železnice. — Lodi.
%  Traduction 2026 d'apres texto/io, controlee sur texto/fr, texto/ru
%  et texto/uk. Consigne : voir texto/cs/10-tabulka-01.tex.
%
%  LE RAIL DONNE ICI LE MEME RESULTAT QUE LA VILLE AU TABLEAU 11, ET
%  POUR LA MEME RAISON HISTORIQUE. Les autres colonnes ont pris le
%  vocabulaire tout fait : « vagón », « perron », « semafor »,
%  « lokomotiva », « tunel », « expres ». Le tcheque l'a pris aussi,
%  mais il en a refait la moitie a dessein, dans les memes annees et
%  par les memes gens : « nádraží » la gare, « kolej » la voie,
%  « výhybka » l'aiguillage, « nástupiště » le quai, « jízdenka » le
%  billet, « zavazadlo » le bagage, « pokladna » le guichet,
%  « průvodčí » le conducteur, « strojvedoucí » le mecanicien.
%
%  ET LE PARTAGE NE SUIT PLUS LA DATE MAIS LE MOT LUI-MEME. « Vagón »
%  est reste, « nádraží » a ete fait ; les deux objets sont arrives le
%  meme jour, par le meme train. Ce qui les separe, c'est qu'un mot
%  d'une syllabe se remplace mal et qu'un mot de trois se remplace
%  bien — et c'est la seule fois du volume ou la LONGUEUR du mot
%  explique quelque chose.
%
%  ZERO INVERSION ATTENDUE, ZERO RELEVEE : ce tableau enumere plus
%  qu'il ne construit, comme dans les six colonnes precedentes.
% ===================================================================

%%K t12-apar-1 apar
\VUsaut{4.0mm}
\VUtitre{15.0pt}{60}{TABULKA č. 12}
\VUsaut{2.4mm}
\VUpk{11.4pt}{40}{Nádraží. --- Železnice. --- Lodi.}
\VUsaut{3.4mm}
\textsc{Poznámka.} --- \textit{Jízdní řády užívané na jednotlivých nádražích jsou různé,
užíváme tedy jako příkladu hodin} \VUgras{francouzské tabulky}.

%%K t12-01-1 p
1. --- Právě jsem cestoval Německem. Před odjezdem jsem si koupil \VUgras{jízdní řád}
\textsuperscript{(15)}, abych zjistil hodiny odjezdu vlaků. Když jsem shledal, že nejvhodnější
vlak odjíždí z Paříže v devět hodin večer a přijíždí do Štrasburku v jednu hodinu třináct
minut, připravil jsem svou cestu a druhého dne jsem se dal dovézt na nádraží.

%%K t12-02-1 p
2. --- Když jsem vstoupil do velké odjezdové haly, za starým venkovským
\VUgras{knězem} \textsuperscript{(2)}, který nesl v ruce svou \VUgras{cestovní tašku}
\textsuperscript{(3)}, mnoho \VUgras{cestujících} \textsuperscript{(1)} už přišlo.

%%K t12-02-2 p
Protože jsem chtěl poslat \VUgras{depeši (telegram)} \textsuperscript{(5)}, dal jsem si od
\VUgras{kontrolora} \textsuperscript{(6)} ukázat \VUgras{telegrafní úřad}
\textsuperscript{(4)}.

%%K t12-03-1 p
3. --- Než jsem přešel do \VUgras{čekárny} \textsuperscript{(7)}, šel jsem si opatřit zpáteční
\VUgras{jízdenku} \textsuperscript{(9)}.

%%K t12-04-1 p
4. --- Nečekal jsem dlouho u \VUgras{okénka} \textsuperscript{(8)}, protože tam bylo málo
lidí. \VUgras{Voják} \textsuperscript{(10)}, který odjížděl na dovolenou, byl tak laskav a
přenechal mi své pořadí, takže jsem měl dost času vyžádat si několik informací od
\VUgras{přednosty stanice} \textsuperscript{(13)}, který rozmlouval s
\VUgras{komisařem} \textsuperscript{(12)} dozoru a se službu konajícím
\VUgras{četníkem} \textsuperscript{(11)}.

%%K t12-tit-1 sub
\VUpk{10.2pt}{40}{Podávání zavazadel.}
\VUsaut{3.0mm}

%%K t12-05-1 p
5. --- Když jsem koupil noviny ve \VUgras{stánku} \textsuperscript{(14)}, dal jsem zvážit svá
\VUgras{zavazadla} \textsuperscript{(17)}, která \VUgras{nosič} \textsuperscript{(16)} právě
přivezl na \VUgras{vozíku na balíky} \textsuperscript{(19)}; \VUgras{úředník}
\textsuperscript{(22)} pověřený
\VUgras{váhou} \textsuperscript{(21)} mi oznámil, že můj kufr váží třicet pět
kilogramů; to činilo \VUgras{nadváhu} pěti kilogramů, za kterou musím doplatit u
\VUgras{okénka} pro zavazadla \textsuperscript{(23)}.

%%K t12-06-1 p
6. --- Protože jsem byl příliš brzy, měl jsem volný čas prohlédnout si pohyb cestujících na
nádraží. Jedni ukládali nebo si vyzvedávali svá lehká zavazadla v
\VUgras{úschovně} \textsuperscript{(25)}; jiní nahlíželi do \VUgras{jízdních řádů}
\textsuperscript{(26)}. Přijíždějící cestující spěchali k \VUgras{východu}
\textsuperscript{(32)} se svým \VUgras{kufříkem} \textsuperscript{(24)} v ruce. Někteří čekali
ve \VUgras{výdejně zavazadel} \textsuperscript{(27)} a předkládali svůj
\VUgras{lístek} \textsuperscript{(29)}, aby si vyzvedli své kufry, své \VUgras{bedny}
\textsuperscript{(28)} nebo své \VUgras{koše} \textsuperscript{(30)}.

%%K t12-07-1 p
7. --- \VUgras{Úředníci potravní daně} \textsuperscript{(33)} pečlivě prohlíželi balíky, aby
zjistili, zda někdo nemá něco k proclení.

%%K t12-08-1 p
8. --- Někteří cestující zašli do \VUgras{bufetu} \textsuperscript{(34)}, kde je
\VUgras{číšník} \textsuperscript{(35)} horlivě obsluhoval. Musí si pospíšit, protože
\VUgras{vlak} \textsuperscript{(36)}, který je rychlík, nezůstane dlouho na nádraží.

%%K t12-09-1 p
9. --- Potom jsem šel na odjezdové nástupiště a hledal jsem přímý \VUgras{vagón}
\textsuperscript{(37)}, abych nebyl nucen měnit vůz během cesty. Nastoupil jsem do vozu s
chodbičkou, který má \VUgras{oddíl} \textsuperscript{(38)} pro kuřáky a oddíl vyhrazený pro
«dámy bez doprovodu».

%%K t12-tit-2 sub
\VUpk{10.2pt}{40}{Odjezd v noci.}
\VUsaut{3.0mm}

%%K t12-10-1 p
10. --- Když jsem vystoupil po \VUgras{stupátku} \textsuperscript{(40)}, zavřel
\VUgras{dveře} \textsuperscript{(39)}, svinul \VUgras{roletu} \textsuperscript{(41)} a
uložil svá malá zavazadla do \VUgras{sítě} \textsuperscript{(43)}, usedl jsem na
\VUgras{lavici} \textsuperscript{(42)}. Když jsem viděl \VUgras{lampáře}
\textsuperscript{(44)} rozsvěcet lampy, pomyslel jsem si, že budu muset strávit noc ve voze, a
zavolal jsem úředníka pověřeného pronájmem \VUgras{podhlavníků} \textsuperscript{(45)}, abych
si o jeden požádal.

%%K t12-11-1 p
11. --- \VUgras{Průvodčí} vlaku přišel zkontrolovat jízdenky. Zeptal jsem se ho, proč jsme
ještě neodjeli, vždyť je správná hodina. Odpověděl mi, že \VUgras{vlakvedoucí}
\textsuperscript{(46)} je zaměstnán ukládáním jízdních kol do \VUgras{služebního vozu}
\textsuperscript{(47)}. Kromě toho ještě nebyly vyměněny všechny ohřívače nohou
\textsuperscript{(48)}.

%%K t12-12-1 p
12. --- Ale brzy jsme měli odjet, protože \VUgras{topič} \textsuperscript{(50)} na svém
\VUgras{tendru} \textsuperscript{(49)} bral uhlí, aby rozdmýchal oheň
\VUgras{lokomotivy} \textsuperscript{(54)}, a \VUgras{strojvedoucí}
\textsuperscript{(51)} čekal jen na znamení \VUgras{zástupce přednosty}
\textsuperscript{(52)}, který byl na \VUgras{nástupišti} \textsuperscript{(53)}.

%%K t12-13-1 p
13. --- \VUgras{Kotel} \textsuperscript{(56)} lokomotivy dutě hučel; \VUgras{komín}
\textsuperscript{(55)} chrlil proudy kouře smíšeného s \VUgras{párou} \textsuperscript{(60)}.
Pronikavá \VUgras{píšťala} \textsuperscript{(59)} mocně zazněla a \VUgras{písty}
\textsuperscript{(58)} se daly do pohybu, pohánějíce \VUgras{kola} \textsuperscript{(57)}
těžkého stroje.

%%K t12-tit-3 sub
\VUsaut{3.0mm}
\VUpk{10.2pt}{40}{Odjezd!}
\VUsaut{3.0mm}

%%K t12-14-1 p
14. --- Rychlost vlaku, běžícího po \VUgras{kolejnicích} \textsuperscript{(64)}, se
zrychlovala; \VUgras{nástupiště} \textsuperscript{(65)} zmizela; minuli jsme
\VUgras{záchodky} \textsuperscript{(66)} a také soupravu vozů odstavenou na druhé
\VUgras{koleji} \textsuperscript{(63)}: \VUgras{lůžkové vozy} \textsuperscript{(67)},
\VUgras{jídelní vůz} \textsuperscript{(68)}, \VUgras{poštovní vůz}
\textsuperscript{(69)}.

%%K t12-noto-1 noto
\VUnotes{143.2mm}{%
\textit{(*) Poznámka.} --- \textit{Rozumí se, že popis cesty je podán podle hranice z doby
před válkou.}%
}

%%K t12-15-1 p
15. --- Potom nám před očima přešlo \VUgras{nákladové nádraží} \textsuperscript{(71)}. Pod
kůlnami nakládali úředníci \VUgras{zboží} \textsuperscript{(72)} odesílané nákladním vlakem.
\VUgras{Posunovači} \textsuperscript{(76)} blízko \VUgras{vodárny} \textsuperscript{(73)}
manévrovali s
\VUgras{dobytčími vozy} \textsuperscript{(75)} a s \VUgras{plošinovými vozy}
\textsuperscript{(79)}, aby sestavili \VUgras{nákladní vlak} \textsuperscript{(80)}.

%%K t12-16-1 p
16. --- Přejeli jsme poslední \VUgras{výhybku} \textsuperscript{(78)}, blízko níž stál
výhybkář; minuli jsme \VUgras{návěstní terč} \textsuperscript{(86)} a nádraží zmizelo v dálce.

%%K t12-17-1 p
17. --- Brzy jsme přejeli \VUgras{železný most} \textsuperscript{(74)}, který vede přes
\VUgras{řeku} \textsuperscript{(81)}. Když jsme přejeli ten most, jeli jsme podél
řeky, po níž mohutné \VUgras{vlečné parníky} \textsuperscript{(82)} táhly těžké
\VUgras{nákladní čluny} \textsuperscript{(83)}. Viděli jsme ještě \VUgras{osobní loď}
\textsuperscript{(84)}, když přistávala u \VUgras{pontonu} \textsuperscript{(85)}.

%%K t12-18-1 p
18. --- Když jsme velmi rychle minuli několik stanic, projeli jsme několika
\VUgras{tunely} \textsuperscript{(87)} a nechali za sebou nesčetné \VUgras{domky}
\textsuperscript{(88)} \VUgras{závoráků.} Ukolébán kymácením vlaku jsem tvrdě usnul.

%%K t12-tit-4 sub
\VUsaut{3.0mm}
\VUpk{10.2pt}{40}{Stanice Deutsch-Avricourt.}
\VUsaut{3.0mm}

%%K t12-19-1 p
19. --- Náhle mě probudil hlas úředníka, který volal: «Deutsch-Avricourt!
\textsuperscript{(*)}. Všichni vystupte!» Konala se celní prohlídka a všechny nudné formality
hranice. Musel jsem srovnat své kapesní hodinky s velkými
\VUgras{hodinami} \textsuperscript{(89)} nádraží, které šly o padesát pět minut
napřed; sotva probuzen jsem stěží rozeznal \VUgras{velkou ručičku} \textsuperscript{(90)} od
\VUgras{malé} \textsuperscript{(91)}; sám
\VUgras{ciferník} \textsuperscript{(92)} byl zcela nezřetelný pro ranní mlhu.

%%K t12-20-1 p
20. --- Zatímco celníci konali svou prohlídku, luštil jsem, zívaje, \VUgras{plakáty}
\textsuperscript{(93)}, které zdobí zdi nádraží. Posnídal jsem, abych ukrátil čas, nahlédl do
mapy německé \VUgras{sítě} \textsuperscript{(94)}, abych viděl, jaká vzdálenost mě ještě dělí
od Štrasburku, a vrátil jsem se do svého vozu posílen nadějí, že brzy dosáhnu cíle své cesty.
